\documentclass[11pt,a4paper]{article}
\usepackage[margin=0.85in,top=1in,bottom=1in]{geometry}
\usepackage{amsmath,amssymb,amsfonts}
\usepackage{graphicx}
\usepackage{float}
\usepackage{enumitem}
\usepackage{microtype}
\usepackage{caption}
\usepackage[hidelinks]{hyperref}
\usepackage[most,breakable]{tcolorbox}
\tcbuselibrary{breakable,skins}
\usepackage[utf8]{inputenc}
\usepackage[T1]{fontenc}
\usepackage{lmodern}
\captionsetup{font=small,labelfont=bf,skip=4pt}
\setlist{leftmargin=1.5em,topsep=2pt}

%% Breakable card — prevents content cutoff across pages
\newtcolorbox{solcard}[3]{
  breakable, enhanced,
  colback=white, colframe=gray!50, arc=2pt,
  title={\textbf{#1}\hfill\textsf{\small #3\,pts}},
  fonttitle=\normalsize\sffamily\bfseries,
  before upper={\small\textit{#2}\par\smallskip\hrule height 0.4pt\smallskip},
  top=4pt, bottom=6pt, left=7pt, right=7pt,
  parbox=false,
}

\title{\textbf{Earnshaw Balls --- Solution}}
\author{\normalsize Y25 (2025) — English translation}
\date{}
\begin{document}\maketitle\vspace{-0.5em}
\begin{solcard}{A1}{Express $\ddot{\vec r}$ and $\dot{\vec\omega}$ in terms of $\vec F$, $m$, $\beta$, $\vec a$.}{0.50}

Let's write down Newton's second law: $$m\ddot{\vec r} = m\vec g + \vec N + \vec F.$$Sphere not move in vertical direction, therefore $m\vec g + \vec N = 0$. We get:


The second answer follows from the law of change in angular momentum relative to the center of the ball:

\par\smallskip\noindent\textbf{Answer:} $$\ddot{\vec r} = \frac{\vec F}{m}$$
\end{solcard}

\begin{solcard}{A2}{Write down the condition for kinematic coupling. Express your answer in terms of $\dot{\vec r}$, $\vec a$, $\vec r$, $\vec \omega$ and $\vec \Omega$.}{0.30}

The velocities of the table and the ball at the point of their contact must match, therefore:

\par\smallskip\noindent\textbf{Answer:} $$\left[\vec\Omega,\vec r\right]=\dot{\vec r} + \left[\vec\omega,\vec a\right]$$
\end{solcard}

\begin{solcard}{A3}{Express $\ddot{\vec r}$ in terms of $\beta$, $\dot{\vec r}$ and $\vec \Omega$.}{0.70}

Let's express $\vec F$ from the equation by $\dot{\vec\omega}$, multiplying it vectorially by $\vec a$: $$\beta m a^2 \left[\dot{\vec \omega},\vec a\right] = \left[ \left[\vec a, \vec F\right],\vec a\right] = \left[\vec a, \left[\vec F, \vec a\right]\right] = a^2\vec F.$$$$\vec F = \beta m\left[\dot{\vec \omega},\vec a\right].$$Express through define quantity, differentiate kinematic relation: $$\left[\vec\Omega,\dot{\vec r}\right]=\ddot{\vec r} + \left[\dot{\vec\omega},\vec a\right].$$$$\vec F = \beta m\left[\dot{\vec \omega},\vec a\right] = \beta m\left(\left[\vec\Omega,\dot{\vec r}\right]-\ddot{\vec r}\right).$$Substitute the resulting expression and get the answer:

\par\smallskip\noindent\textbf{Answer:} $$\ddot{\vec r} = \frac{\beta}{1+\beta} \left[\vec\Omega, \dot{\vec r}\right]$$
\end{solcard}

\begin{solcard}{A4}{Let the initial position and speed of the ball be $\vec r_0$ and $\vec v_0$, respectively. Determine how $\dot{\vec r}$ depends on $\vec r$ as it moves further. Qualitatively depict the trajectory of the center of the ball in this case and indicate the characteristic dimensions and coordinates.}{0.70}

Note that the expression obtained in the previous paragraph is the equality of two total derivatives with respect to time: $$\frac{\text{d}}{\text{d}t}\dot{\vec r} =\frac{\text{d}}{\text{d}t} \frac{\beta}{1+\beta} \left[\vec\Omega, \vec r\right].$$Integrate this relation taking into account that, that velocity center sphere always parallel plane table: $$\dot{\vec r} = \frac{\beta}{1+\beta} \left[\vec\Omega, \vec r - \vec R_0\right],$$where $\vec R_0$ - constant vector. Note, that this equation define motion by circle with center $\vec R_0$. Now let's substitute the initial conditions. Then integrating the acceleration will give us:


As we have already seen, the trajectory is a circle. All that remains is to find its center and radius. $$\dot{\vec r} = \frac{\beta}{1+\beta} \left[\vec\Omega, \vec r - \vec R_0\right] = \frac{\beta}{1+\beta} \left[\vec\Omega, \vec r - \vec r_0\right] + \vec v_0$$Multiply by vector on $\vec\Omega$: $$\frac{\beta}{1+\beta} \left[\vec\Omega,\left[\vec\Omega, \vec r - \vec R_0\right]\right] = \frac{\beta}{1+\beta} \left[\vec\Omega,\left[\vec\Omega, \vec r - \vec r_0\right]\right] + \left[\vec\Omega,\vec v_0\right].$$$$\frac{\beta}{1+\beta} \left[\vec\Omega,\left[\vec\Omega, \vec r_0 - \vec R_0\right]\right] = \frac{\beta\Omega^2}{1+\beta}\left(\vec R_0 - \vec r_0\right)= \left[\vec\Omega,\vec v_0\right]$$$$\vec R_0 - \vec r_0=\frac{1+\beta}{\Omega^2\beta} \left[\vec\Omega,\vec v_0\right]$$From obtain relation immediately obtain center and radius circle: $$\vec R_0 = \vec r_0+\frac{1+\beta}{\Omega^2\beta} \left[\vec\Omega,\vec v_0\right],\qquad \rho = \frac{\left(1+\beta\right)v_0}{\Omega\beta}.$$

\par\smallskip\noindent\textbf{Answer:} $$\dot{\vec r} = \frac{\beta}{1+\beta} \left[\vec\Omega, \vec r - \vec r_0\right] + \vec v_0$$
\begin{figure}[H]
  \centering
  \includegraphics[width=0.72\textwidth]{T3_S_EN_assets/fig1.jpg}
\end{figure}
\end{solcard}

\begin{solcard}{B1}{Express $\ddot{\vec r}$ in terms of $\vec g_\parallel$, $\beta$, $\dot{\vec r}$ and $\vec \Omega$.}{0.30}

The only difference from the previous part is that in the initial expression for $\ddot{\vec r}$ there will be a non-zero contribution $m\vec g + \vec N = m\vec g_\parallel$. $$\ddot{\vec r} = \vec g_\parallel+\beta\left(\left[\vec\Omega,\dot{\vec r}\right]-\ddot{\vec r}\right)$$

\par\smallskip\noindent\textbf{Answer:} $$ \ddot{\vec r} = \frac{\vec g_\parallel}{1+\beta}+\frac{\beta}{1+\beta}\left[\vec\Omega,\dot{\vec r}\right] $$
\end{solcard}

\begin{solcard}{B2}{Let us denote $$\vec{u} = \frac{\left[\vec{\Omega} \times \vec g_\parallel\right]}{\Omega^2}.$$ Qualitatively depict trajectory center sphere for following initial velocity: $\vec{v}_0 = \vec{u}$ $\vec{v}_0 = 2.5 \vec{u}$ $\vec{v}_0 = 5 \vec{u}$ $\vec{v}_0 = 0$ For case 4 calculate numerically characteristic size trajectory for $g_\parallel=1.0~m/s^2$, $\Omega = 10~rad/s$.}{1.40}

For case 4, numerically calculate the characteristic dimensions of the trajectory for $g_\parallel=1.0~m/s^2$, $\Omega = 10~\mathrm{rad}/with$.


Let's analyze the answer from the previous paragraph. Acceleration consists of a constant term and a term obtained by the vector product of speed and a constant vector. This is completely analogous to motion in a crossed magnetic and electric field. Let's find the drift speed $\dot{\vec r} = \vec v + \vec v_\text{d}$. $$ \dot{\vec v} = \frac{\vec g_\parallel}{1+\beta}+\frac{\beta}{1+\beta}\left[\vec\Omega,\vec v\right]+\frac{\beta}{1+\beta}\left[\vec\Omega,\vec v_\text{d}\right] $$$$ 0 = \frac{\vec g_\parallel}{1+\beta}+\frac{\beta}{1+\beta}\left[\vec\Omega,\vec v_\text{d}\right] $$Vector multiply by on $\vec\Omega$, obtain value velocity drift: $$ \vec v_\text{d} = \frac{\left[\vec \Omega, \vec g_\parallel\right]}{\Omega^2\beta} = \frac{5}{2}\vec u. $$Vector $\vec v$ rotate with constant angular velocity $\beta \vec\Omega/(1+\beta)$. Using the obtained expressions, we obtain the form of the trajectory for each case.

\par\smallskip\noindent\textbf{Answer:} Case 1: the difference from the drift speed is small, so there will be a trajectory without self-intersections.
\begin{figure}[H]
  \centering
  \includegraphics[width=0.72\textwidth]{T3_S_EN_assets/fig2.jpg}
\end{figure}
\begin{figure}[H]
  \centering
  \includegraphics[width=0.72\textwidth]{T3_S_EN_assets/fig3.jpg}
\end{figure}
\begin{figure}[H]
  \centering
  \includegraphics[width=0.72\textwidth]{T3_S_EN_assets/fig4.jpg}
\end{figure}
\begin{figure}[H]
  \centering
  \includegraphics[width=0.72\textwidth]{T3_S_EN_assets/fig5.jpg}
\end{figure}
\end{solcard}

\begin{solcard}{C1}{Express $\ddot{\vec r}$ in terms of $m$, $\beta$, $\vec a$, $\vec r$, $\vec N$, $\vec g$, $\dot{\vec \omega}$}{0.20}

From Newton's second law: $$ \ddot{\vec r} = \frac{\vec N}{m} + \vec g + \frac{\vec F}{m}. $$From the law change moment momentum: $$ m a^2 \beta \dot{\vec \omega} = \left[\vec a \times \vec F\right] \quad \Rightarrow \quad ma^2 \beta \left[\vec a \times \dot{\vec \omega}\right] = \left[\vec a \times \left[\vec a \times \vec F\right]\right] = \vec a \left(\vec a \cdot \vec F\right) - \vec F a^2 = - \vec F a^2. $$From:

\par\smallskip\noindent\textbf{Answer:} $$ \ddot{\vec r} = \frac{\vec N}{m} + \vec g - \left[\vec a \times \dot{\vec \omega}\right] \beta $$
\end{solcard}

\begin{solcard}{C2}{Write down the condition for kinematic coupling. Express your answer in terms of $\dot{\vec r}$, $\vec a$, $\vec r$, $\vec \omega$, $\vec \Omega$.}{0.40}

The velocities of the table and the ball at the point of their contact must coincide, therefore: $$ \dot{\vec r} + \left[\vec \omega \times \vec a\right] = \left[\vec \Omega \times \left(\vec r + \vec a\right)\right], $$from:

\par\smallskip\noindent\textbf{Answer:} $$ \dot{\vec r} = \left[\vec \Omega \times \left(\vec r + \vec a\right)\right] - \left[\vec \omega \times \vec a\right] $$
\end{solcard}

\begin{solcard}{C3}{Express $\ddot{\vec r}$ in terms of $m$, $\beta$, $\vec a$, $\vec N$, $\vec g$, $\dot{\vec r}$, $\vec\Omega$, $\vec \zeta$, $\vec \omega$.}{0.90}

Let us differentiate the expression for $\dot{\vec r}$ with respect to time: $$ \ddot{\vec r} = \left[\vec \Omega \times \dot{\vec r}\right] + \left[\vec \Omega \times \dot{\vec a}\right] - [\dot{\vec \omega} \times \vec a] - [\vec \omega \times \dot{\vec a}] $$ From the expression for $\ddot{\vec r}$, obtain in part C1 : $$ [\dot{\vec \omega} \times \vec a] = \frac{1}{\beta} \left(\ddot{\vec r} - \frac{\vec N}{m} - \vec g\right) $$ Also: $$ \dot{\vec a} = \left[\vec \zeta \times \vec a\right], $$ from where:


Let's see how this formula is transformed taking into account some approximations, namely: $\theta \ll 1$, that is, the table is slightly different from a flat one and the influence of gravity is weak; the motion of the ball is finite, that is, it moves in a limited region of space near the vertex and does not move away to infinity; $a \lesssim r$, that is, the ball is not too close to the top of the cone; $\zeta \lesssim \Omega$, that is, the rotation of the table occurs quite quickly in comparison with the speed of the ball; $\left|\left(\vec\omega, \vec a\right) \right|\lesssim \Omega r$, that is, the spinning speed is not too high. Let us restrict ourselves to terms of order of smallness $\theta$. Note that the terms with $\vec N$ and $\vec g$ have the order of smallness $\theta$, since the projection $\vec g$ onto the surface of the cone is of this order. Due to the fact that the motion is finite, the remaining terms must be of the same order. Now it remains to show that the first term is an order of magnitude larger than the last, then the first is of the first order of smallness, and the last is of the second, and it can be discarded. Estimation of the value of the first term: $$\left|\beta\frac{[\vec \Omega, \dot{\vec r}]}{1 + \beta}\right|\sim \left|[\vec \Omega, \dot{\vec r}]\right|\sim \Omega \dot r + \Omega \zeta r.$$ Now last term expand on two term and let's analyze each from they taking into account that, that angle between $\vec\zeta$ and $\vec a$ equal $\theta$: $$\left|\frac{\beta}{1 + \beta}\left[\vec \Omega, \left[\vec \zeta, \vec a\right]\right]\right|\sim\Omega\zeta a\theta \lesssim \Omega\zeta r\theta \ll \Omega\zeta r \lesssim \left|\beta\frac{[\vec \Omega, \dot{\vec r}]}{1 + \beta}\right|.$$ For estimate second term expand $\vec\omega$ on component along and perpendicularly $\vec a$: $\vec \omega = \vec \omega_a + \vec \omega_\perp$. Note, that from kinetic relation and last condition approximation: $$\omega a \lesssim \left|\dot{\vec r}\right|+\Omega r +\left|\left(\vec\omega, \vec a\right) \right|\sim \dot r + \Omega r.$$ Thus last expression, smallness which must prove: $$\left|\frac{\beta}{1 + \beta}\left[\vec \omega, \left[\vec \zeta, \vec a\right]\right]\right|\sim\omega\zeta a\theta \lesssim \left(\dot r + \Omega r\right)\zeta \theta \ll \Omega \dot r + \Omega \zeta r\sim \left|\beta\frac{[\vec \Omega, \dot{\vec r}]}{1 + \beta}\right|.$$ Thus, we obtain the desired approximate formula.


Let us restrict ourselves to terms of order of smallness $\theta$.


Note that the terms with $\vec N$ and $\vec g$ have the order of smallness $\theta$, since the projection $\vec g$ onto the surface of the cone is of this order. Due to the fact that the motion is finite, the remaining terms must be of the same order. Now it remains to show that the first term is an order of magnitude larger than the last, then the first is of the first order of smallness, and the last is of the second, and it can be discarded.


Estimation of the value of the first term:


$$\left|\beta\frac{[\vec \Omega, \dot{\vec r}]}{1 + \beta}\right|\sim \left|[\vec \Omega, \dot{\vec r}]\right|\sim \Omega \dot r + \Omega \zeta r.$$


Now let’s decompose the last term into two terms and analyze each of them, taking into account the fact that the angle between $\vec\zeta$ and $\vec a$ is equal to $\theta$:


$$\left|\frac{\beta}{1 + \beta}\left[\vec \Omega, \left[\vec \zeta, \vec a\right]\right]\right|\sim\Omega\zeta a\theta \lesssim \Omega\zeta r\theta \ll \Omega\zeta r \lesssim \left|\beta\frac{[\vec \Omega, \dot{\vec r}]}{1 + \beta}\right|.$$


To estimate the second term, we decompose $\vec\omega$ into components along and perpendicular to $\vec a$: $\vec \omega = \vec \omega_a + \vec \omega_\perp$. Note that from the kinetic connection and the last approximation condition:


$$\omega a \lesssim \left|\dot{\vec r}\right|+\Omega r +\left|\left(\vec\omega, \vec a\right) \right|\sim \dot r + \Omega r.$$


Thus, the last expression, the smallness of which must be proven:


$$\left|\frac{\beta}{1 + \beta}\left[\vec \omega, \left[\vec \zeta, \vec a\right]\right]\right|\sim\omega\zeta a\theta \lesssim \left(\dot r + \Omega r\right)\zeta \theta \ll \Omega \dot r + \Omega \zeta r\sim \left|\beta\frac{[\vec \Omega, \dot{\vec r}]}{1 + \beta}\right|.$$


Thus we obtain the desired approximate formula.

\par\smallskip\noindent\textbf{Answer:} $$ \ddot{\vec r} = \frac{[\vec \Omega \times \dot{\vec r}]}{1 + \beta^{-1}} + \frac{1}{1 + \beta} \left( \frac{\vec N}{m} + \vec g\right) + \frac{1}{1 + \beta^{-1}} \left[\left(\vec \Omega - \vec \omega\right) \times \left[\vec \zeta \times \vec a\right]\right] $$
\end{solcard}

\begin{solcard}{C4}{Prove that in the considered approximation the quantity $$ l = r^2 \left( \zeta-\cfrac{\Omega\beta}{2\left(1+\beta\right)} \right). $$Here $\zeta = |\vec{\zeta}|$ is conserved.}{1.00}

First decision. Let us expand the terms from this equation according to the introduced basis: $$ \dot{\vec r} = \left[\vec \zeta \times \vec r\right] + \dot r \vec e_r = \zeta r \vec e_\varphi + \dot r \vec e_r $$ $$ \ddot{\vec r} = \left( \dot \zeta r \vec e_\varphi + \zeta \dot r \vec e_\varphi + \zeta r \frac{\text{d} \vec e_\varphi}{\text{d} t} \right) + \left( \ddot r \vec e_r + \dot r \frac{\text{d} \vec e_r}{\text{d} t} \right) =\\= \dot \zeta r \vec e_\varphi + \zeta \dot r \vec e_\varphi + \zeta^2 r \left(-\vec e_r - \theta \vec e_z\right) + \ddot r \vec e_r + \zeta \dot r \vec e_\varphi = \left(\dot \zeta r + 2 \zeta \dot r\right) \vec e_\varphi + \left(\ddot r - \zeta^2 r\right) \vec e_r + \zeta^2 r \theta \vec e_z $$ $$ \left[\vec \Omega \times \dot{\vec r}\right] = \Omega \dot r \vec e_\varphi - \Omega \zeta r \vec e_r + \Omega \zeta r \theta \vec e_z $$ $$ \frac{\vec N}{m} + \vec g = g \theta \vec e_r + \left( \frac{N_z}{m} - g \right) \vec e_z $$ Now substitute they in given equation, take projection on $\vec e_\varphi$: $$ \dot \zeta r + 2 \zeta \dot r = \frac{\Omega \dot r}{1 + \beta^{-1}} \quad \Rightarrow \quad \frac{\text d \zeta}{\text d r} = \frac{1}{r} \left( \frac{\Omega}{1 + \beta^{-1}} - 2 \zeta \right) \quad \Rightarrow\\ \int \frac{\text d \zeta}{\dfrac{\Omega}{1 + \beta^{-1}} - 2 \zeta} = \int \frac{\text d r}{r} \quad \Rightarrow \quad r^2 \left( \zeta-\cfrac{\Omega\beta}{2\left(1+\beta\right)} \right) = \text{const}, $$ which is what needed to be proved.


Let us expand the terms from this equation according to the introduced basis: $$ \dot{\vec r} = \left[\vec \zeta \times \vec r\right] + \dot r \vec e_r = \zeta r \vec e_\varphi + \dot r \vec e_r $$ $$ \ddot{\vec r} = \left( \dot \zeta r \vec e_\varphi + \zeta \dot r \vec e_\varphi + \zeta r \frac{\text{d} \vec e_\varphi}{\text{d} t} \right) + \left( \ddot r \vec e_r + \dot r \frac{\text{d} \vec e_r}{\text{d} t} \right) =\\= \dot \zeta r \vec e_\varphi + \zeta \dot r \vec e_\varphi + \zeta^2 r \left(-\vec e_r - \theta \vec e_z\right) + \ddot r \vec e_r + \zeta \dot r \vec e_\varphi = \left(\dot \zeta r + 2 \zeta \dot r\right) \vec e_\varphi + \left(\ddot r - \zeta^2 r\right) \vec e_r + \zeta^2 r \theta \vec e_z $$ $$ \left[\vec \Omega \times \dot{\vec r}\right] = \Omega \dot r \vec e_\varphi - \Omega \zeta r \vec e_r + \Omega \zeta r \theta \vec e_z $$ $$ \frac{\vec N}{m} + \vec g = g \theta \vec e_r + \left( \frac{N_z}{m} - g \right) \vec e_z $$ Now substitute they in given equation, take projection on $\vec e_\varphi$: $$ \dot \zeta r + 2 \zeta \dot r = \frac{\Omega \dot r}{1 + \beta^{-1}} \quad \Rightarrow \quad \frac{\text d \zeta}{\text d r} = \frac{1}{r} \left( \frac{\Omega}{1 + \beta^{-1}} - 2 \zeta \right) \quad \Rightarrow\\ \int \frac{\text d \zeta}{\dfrac{\Omega}{1 + \beta^{-1}} - 2 \zeta} = \int \frac{\text d r}{r} \quad \Rightarrow \quad r^2 \left( \zeta-\cfrac{\Omega\beta}{2\left(1+\beta\right)} \right) = \text{const}, $$ which is what needed to be proved.


Second solution. Let's consider the change in the angular momentum of the center of mass: $$\frac{\text{d}}{\text{d}t}\left[\vec r, \dot{\vec r}\right] = \left[\vec r, \ddot{\vec r}\right].$$ For projection on axis $z$: $$\frac{\text{d}}{\text{d}t}\left(\vec e_z, \vec r, \dot{\vec r}\right) = \frac{\text{d}}{\text{d}t}\left(r^2\zeta\right) = \left(\vec e_z, \vec r, \ddot{\vec r}\right).$$ Note, that $\left(\vec e_z, \vec r, \vec N\right)$ and $\left(\vec e_z, \vec r, \vec g\right)$ equal zero, since this triple vector lie in one plane, therefore from acceleration remain only one term: $$\frac{\text{d}}{\text{d}t}\left(r^2\zeta\right) = \frac{\beta}{1+\beta} \left(\vec e_z, \vec r,  [\vec \Omega, \dot{\vec r}]\right) = \frac{\beta}{1+\beta} \left(\vec e_z, \vec\Omega\left(\vec r,\dot{\vec r}\right) - \dot{\vec r} \left(\vec\Omega, \vec r\right)\right).$$ $$\frac{\text{d}}{\text{d}t}\left(r^2\zeta\right) = \frac{\beta}{1+\beta} \left(\left(\vec e_z, \vec\Omega\right)\left(\vec r,\dot{\vec r}\right) - \left(\vec e_z,\dot{\vec r}\right) \left(\vec\Omega, \vec r\right)\right)$$ Substitute the value of the scalar products and discard term second order smallness: $$\frac{\text{d}}{\text{d}t}\left(r^2\zeta\right) = \frac{\beta}{1+\beta} \left(\Omega\frac{\text{d}}{\text{d}t}\left(\frac{r^2}{2}\right) - r\dot{r}\Omega \theta^2\right) \approx \frac{\beta\Omega}{1+\beta}\frac{\text{d}}{\text{d}t}\left(\frac{r^2}{2}\right) .$$ $$\frac{\text{d}}{\text{d}t}\left(r^2 \left( \zeta-\cfrac{\Omega\beta}{2\left(1+\beta\right)} \right)\right) = 0$$ Obtain, that quantity $l$ is saved.


Let's consider the change in the angular momentum of the center of mass:


$$\frac{\text{d}}{\text{d}t}\left[\vec r, \dot{\vec r}\right] = \left[\vec r, \ddot{\vec r}\right].$$


For projection onto the $z$ axis:


$$\frac{\text{d}}{\text{d}t}\left(\vec e_z, \vec r, \dot{\vec r}\right) = \frac{\text{d}}{\text{d}t}\left(r^2\zeta\right) = \left(\vec e_z, \vec r, \ddot{\vec r}\right).$$


Note that $\left(\vec e_z, \vec r, \vec N\right)$ and $\left(\vec e_z, \vec r, \vec g\right)$ are equal to zero, since these triples of vectors lie in the same plane, so only one term will remain from the acceleration:


$$\frac{\text{d}}{\text{d}t}\left(r^2\zeta\right) = \frac{\beta}{1+\beta} \left(\vec e_z, \vec r,  [\vec \Omega, \dot{\vec r}]\right) = \frac{\beta}{1+\beta} \left(\vec e_z, \vec\Omega\left(\vec r,\dot{\vec r}\right) - \dot{\vec r} \left(\vec\Omega, \vec r\right)\right).$$


$$\frac{\text{d}}{\text{d}t}\left(r^2\zeta\right) = \frac{\beta}{1+\beta} \left(\left(\vec e_z, \vec\Omega\right)\left(\vec r,\dot{\vec r}\right) - \left(\vec e_z,\dot{\vec r}\right) \left(\vec\Omega, \vec r\right)\right)$$


Let us substitute the values of the rock products and discard the terms of the second order of smallness:


$$\frac{\text{d}}{\text{d}t}\left(r^2\zeta\right) = \frac{\beta}{1+\beta} \left(\Omega\frac{\text{d}}{\text{d}t}\left(\frac{r^2}{2}\right) - r\dot{r}\Omega \theta^2\right) \approx \frac{\beta\Omega}{1+\beta}\frac{\text{d}}{\text{d}t}\left(\frac{r^2}{2}\right) .$$


$$\frac{\text{d}}{\text{d}t}\left(r^2 \left( \zeta-\cfrac{\Omega\beta}{2\left(1+\beta\right)} \right)\right) = 0$$


We found that the value $l$ is preserved.

\end{solcard}

\begin{solcard}{C5}{Express $\ddot{r}$ in terms of $r$, $\zeta$, $\beta$, $\Omega$, $\theta$ and $g$.}{0.50}

Now consider the equation: $$ \ddot{\vec r} = \frac{1}{1+\beta} \left(\vec g + \frac{\vec N}{m}\right) + \frac{\beta\left[\vec \Omega \times \dot{\vec r}\right]}{1 + \beta}. $$Substitute in it projection term on $\vec e_r$: $$ \ddot r - \zeta^2 r = \frac{g \theta}{1 + \beta} - \frac{\beta \Omega \zeta r}{1 + \beta} \quad \Rightarrow \quad \ddot r = \zeta^2 r + \frac{g \theta}{1 + \beta} - \frac{\beta \Omega \zeta r}{1 + \beta}. $$

\par\smallskip\noindent\textbf{Answer:} $$ \ddot r = \zeta^2 r + \frac{g \theta}{1 + \beta} - \frac{\beta \Omega \zeta r}{1 + \beta} $$
\end{solcard}

\begin{solcard}{C6}{Determine at what values of $r$ motion along circular trajectories around the axis of symmetry of the cone $(r = \text{const})$ is possible and determine the values of $\zeta$ corresponding to such motion. Express your answer in terms of $r$, $\beta$, $\Omega$, $\theta$ and $g$.}{0.70}

$$ r = \text{const} \quad \Rightarrow \quad \ddot r = 0 = \zeta^2 r + \frac{g \theta}{1 + \beta} - \frac{\beta \Omega \zeta r}{1 + \beta}. $$Then obtain square equation with respect to $\zeta$: $$ \zeta^2 - \zeta \cdot \frac{\Omega}{1 + \beta^{-1}} + \frac{g \theta}{(1 + \beta) r} = 0 \quad \Rightarrow \quad \zeta_{1,2} = \frac{\Omega}{2(1 + \beta^{-1})} \pm \sqrt{\left(\frac{\Omega}{2(1 + \beta^{-1})}\right)^2 - \frac{g \theta}{(1 + \beta) r}}. $$Since the expression under the radical is non-negative: $$ r \ge \frac{4 g \theta (1 + \beta)}{\Omega^2 \beta^2} = r_{\min}. $$

\par\smallskip\noindent\textbf{Answer:} $$ \zeta_{1,2} = \frac{\Omega}{2(1 + \beta^{-1})} \pm \sqrt{\left(\frac{\Omega}{2(1 + \beta^{-1})}\right)^2 - \frac{g \theta}{(1 + \beta) r}} $$$$ r \ge \frac{4 g \theta (1 + \beta)}{\Omega^2 \beta^2} $$
\end{solcard}

\begin{solcard}{C7}{Express $\dot{r}^2$ up to an arbitrary constant in terms of $l$, $r$, $g$, $\beta$, $\Omega$ and $\theta$. Qualitatively depict the appearance of phase diagrams $\dot{r}(r)$.}{0.80}

It is easy to see that: $$ \ddot r = \frac{\text d \dot r}{\text d t} = \frac{\text d \dot r}{\text d r / \dot r} = \frac{\text d ({\dot r}^2)}{2 \text d r} \quad \Rightarrow \quad {\dot r}^2 = 2 \int \ddot r(r) \text d r. $$We know, that be conserved quantity $l$, consequently: $$ \zeta (r) = \frac{\Omega\beta}{2\left(1+\beta\right)} + \frac{l}{r^2}. $$Substitute $\zeta$ in expression for $\ddot r$, obtain: $$ {\dot r}^2 = 2 \int \left( \frac{g \theta}{1 + \beta} + \frac{l^2}{r^3} - \frac{\Omega^2 r}{4 (1 + \beta^{-1})^2} \right) \text d r = \frac{2 g \theta r}{1 + \beta} - \frac{\Omega^2 r^2}{4 (1 + \beta^{-1})^2} - \frac{l^2}{r^2} + \text{const}. $$

\par\smallskip\noindent\textbf{Answer:} $$ {\dot r}^2 = \frac{2 g \theta r}{1 + \beta} - \frac{\Omega^2 r^2}{4 (1 + \beta^{-1})^2} - \frac{l^2}{r^2} + \text{const} $$
\begin{figure}[H]
  \centering
  \includegraphics[width=0.72\textwidth]{T3_S_EN_assets/fig6.jpg}
\end{figure}
\end{solcard}

\begin{solcard}{C8}{Let us consider a small perturbation of a circular orbit with radius $r_0$. Determine the cyclic frequency $\xi$ of small radial vibrations. Express your answer in terms of $\Omega$, $r_0$, $\theta$, $g$, $\beta$.}{0.40}

The orbit was disturbed so that the distance from the axis to the center of the ball was presented in the form: $$ r(t) = r_0 + \delta (t), ~(\text{where } \delta \ll r_0 \text{ and } \left< r \right> = r_0). $$Then: $$ \ddot \delta = - \xi^2 \delta = \ddot r(r_0) = \frac{2g\theta}{1 + \beta} + \frac{l^2}{r^3} - \frac{\Omega^2 r}{4 (1 + \beta^{-1})^2} \approx - \frac{3 l^2 \delta}{r_0^4} - \frac{\Omega^2 \delta}{4 (1 + \beta^{-1})^2}, $$from which: $$ \xi = \sqrt{\frac{3 l^2}{r_0^4} + \frac{\Omega^2}{4 (1 + \beta^{-1})^2}}, $$where: $$ l = r_0^2 \left(\zeta (r_0) - \frac{\Omega}{2 (1 + \beta^{-1})} \right) = \pm r_0^2 \sqrt{\frac{\Omega^2}{4 (1 + \beta^{-1})^2} - \frac{g \theta}{(1 + \beta) r_0}}. $$Substitute $l$ in expression for $\xi$, we get:

\par\smallskip\noindent\textbf{Answer:} $$ \xi = \sqrt{\frac{\Omega^2}{(1 + \beta^{-1})^2} - \frac{3 g \theta}{(1 + \beta) r_0}} $$
\end{solcard}

\begin{solcard}{C9}{Let the ball initially move with $\dot{\varphi}=\zeta_0=0.44~\mathrm{rad}/with$, $\dot{r} = 0$, $r_0 = 15~cm$ along a cone with $\theta = 1^\circ$ and $\Omega = 10~\mathrm{rad}/with$. Draw the type of trajectory and indicate its characteristic dimensions. It is known that $g = 9.81 ~m/s^2$. Note: you can consider it known that the deviation from the circular orbit in this case is small, therefore the dependence of $\ddot{r}(r)$ and $\dot\varphi (r)$ can be approximated by a linear function near the AVERAGE value $r$.}{1.20}

Note: you can consider it known that the deviation from the circular orbit in this case is small, therefore the dependence of $\ddot{r}(r)$ and $\dot\varphi (r)$ can be approximated by a linear function near the AVERAGE value $r$.


Let's find the value of $l$ using these values: $$l = -0.022243~\text{m}^2/\text{s}.$$in initial moment time $\ddot r$: $$\ddot r(0) = \frac{2g\theta}{1 + \beta} + \frac{l^2}{r_0^3} - \frac{\beta^2\Omega^2 r_0}{4 (1 + \beta)^2} = -0.037232~m/s^2.$$We obtain, that in initial moment time $r$ maximally. Find amplitude oscillation $A$. Note, that mean value radius-vector be $r_0-A$. $$0.037232~m/s^2 = \xi^2 A$$$$\xi^2 = \frac{3 l^2}{\left(r_0-A\right)^4} + \frac{\beta^2\Omega^2}{4 (1 + \beta)^2}$$We obtain an equation for $A$: $$0.037232~\text{m}/\text{s}^2 = \left(\frac{3 l^2}{(r_0-A)^4} + \frac{\beta^2\Omega^2}{4(1+\beta)^2}\right) A.$$ Note that in right part monotonic by $A$ function, therefore solution uniquely. Find it numerically: $$A = 6.6961~\mathrm{mm}.$$Find $\xi$ and mean value $\zeta$: $$\xi = \sqrt{\frac{3 l^2}{\left(r_0-A\right)^4} + \frac{\beta^2\Omega^2}{4 (1 + \beta)^2}}=2.3580~\text{s}^{-1},$$$$\left<\zeta\right> = \frac{\Omega\beta}{2\left(1+\beta\right)} + \frac{l}{(r_0-A)^2}=0.34545~\mathrm{rad}/\text{s}.$$Then equations of motion: $$\begin{cases} r = r_0 - A + A\cos\left(\xi t\right), \\~\\ \zeta = \left<\zeta\right> + \left(\zeta_0-\left<\zeta\right>\right)\cos\left(\xi t\right). \end{cases}$$It is worth noting that it is extremely important to take into account that the mean value of $\zeta$ does not coincide with the initial one, since the difference of these values is very significant. In this case the harmonicity in time of $\zeta$ and $r$ holds, because $r$ varies weakly and the linear approximations of all equations of motion are reasonable. The qualitative form of the trajectory is already clear: motion with periodic change of the radius, on average the change of the angle $\varphi$ occurs with velocity $ \left<\zeta\right>$. Number of radial oscillations per one revolution: $$N = \frac{\xi}{\left<\zeta\right>} = 6.83.$$It is worth noting, that in the numerical solution of the differential equations of motion are solved, it is obtained, that in fact $N=7$ and trajectory is closed with good accuracy, the approximation gives a close answer, that confirms the reasonableness of the approximations made. Let us draw the trajectory we obtained: $$\begin{cases} r = r_0 - A + A\cos\left(\xi t\right), \\~\\ \varphi = \left<\zeta\right>t + \dfrac{\zeta_0-\left<\zeta\right>}{\xi}\sin\left(\xi t\right). \end{cases}$$It lies between circles with radii $r_0 = 15.0~cm$ and $r_0 - 2A = 13.7~cm$ (they are shown in the figure below with a dotted line) and almost 7 radial oscillations occur during one revolution along the trajectory:

\begin{figure}[H]
  \centering
  \includegraphics[width=0.72\textwidth]{T3_S_EN_assets/fig7.jpg}
\end{figure}
\end{solcard}

\end{document}